\documentclass{utmabstract}

% The official abstract template does not require a logo. Enable only if requested.
% \UTMShowLogotrue
\utmlogo{../../shared/assets/UTM_Logo.png}

\abstracttitle{Titlul lucrării}
\abstractauthors{Prenumele NUMELE autor unu$^{1}$, Prenumele NUMELE autor doi$^{2*}$, Prenumele NUMELE autor trei$^{3}$}
\abstractaffiliations{$^{1}$Departamentul, grupa, Facultatea, Universitatea Tehnică a Moldovei, Chișinău, Republica Moldova\\
$^{2}$Departamentul, grupa, Facultatea, Universitatea Tehnică a Moldovei, Chișinău, Republica Moldova\\
$^{3}$Departamentul, grupa, Facultatea, Universitatea Tehnică a Moldovei, Chișinău, Republica Moldova}
\abstractcorresponding{*Autorul corespondent: Prenumele Numele, email@utm.md}
\abstractsupervisor{\textbf{Îndrumătorul/coordonatorul științific Prenumele NUMELE}, titlul științific, afilierea}

\begin{document}
\makeutmabstracttitle

\begin{rezumat}
Rezumatul trebuie să fie clar, concis și concentrat pe subiect. Textul va prezenta scopul cercetării, metodologia utilizată, principalele rezultate obținute și concluziile relevante. Se recomandă evitarea afirmațiilor generale și evidențierea elementului de noutate al lucrării, astfel încât cititorul să înțeleagă rapid contribuția autorilor față de literatura de specialitate. Formulările trebuie să fie precise, iar prescurtările trebuie explicate la prima apariție. Pentru lucrările redactate în limba română se folosesc diacritice. Rezumatul se redactează cu Times New Roman, 12 pt, italic, aliniere justificată și are aproximativ o jumătate de pagină.
\end{rezumat}

\cuvintecheie{modelare, optimizare, analiză, validare, rezultate}

\end{document}
